% from ega2-1-fr.tex
\oldpage[II]{123}

remplacer $A$ par $A[T]$ (qui est noethérien) et $E$ par
$E\otimes_AA[T]=E'$ qui est tel que
$E'\otimes_{A[T]}K[T]=E\otimes_AK[T]$ soit un $K[T]$-module libre de rang $n$.
Revenant aux notations initiales, on voit donc qu'il suffit de prouver que le
$A$-module engendré par les éléments $\det u$, où $u$ parcourt
$\operatorname{Hom}_A(E,E)$, est de type fini~; \emph{a fortiori} (puisque
tout sous-module d'un $A$-module de type fini est de type fini) il suffit de
prouver que lorsque $v$ parcourt l'ensemble des endomorphismes de
$\bigwedge^nE$, le $A$-module engendré par les $\det v$ est de type fini~;
autrement dit, on est encore ramené au cas où $n=1$. Mais alors la proposition
résulte de ce que $\operatorname{Hom}_A(E,E)$ est un $A$-module de type fini
et de ce que $v\mapsto\det v$ est un homomorphisme de ce $A$-module dans $K$.

Soit $F$ le noyau de l'homomorphisme canonique $E\to E\otimes_AK$ et soit
$E_0=E/F$~; on voit comme ci-dessus que $E\otimes_AK$ s'identifie à
$E_0\otimes_AK$, que $u(F)\subset F$, et que si $u_0$ est l'endomorphisme de
$E_0$ déduit de $u$ par passage aux quotients, $u\otimes1$ s'identifie à
$u_0\otimes1$, et $\sigma_j(u)=\sigma_j(u_0)$ pour tout $j$. \emph{Si l'on a
$F=0$}, les formules (\hyperref[II.6.4.1.3-fr]{6.4.1.3}) et
(\hyperref[II.6.4.1.5-fr]{6.4.1.5}) ont un sens et sont valables lorsque $E$
est identifié à un sous-module de $E\otimes_AK$ et les $u^j$ à des
homomorphismes $E\to E\otimes_AK$~; par suite, la prop.
(\hyperref[II.6.4.5-fr]{6.4.5}) s'étend aussi à ce cas, avec sa démonstration.
\end{remark}

\begin{env}[6.4.8]
\label{II.6.4.8-fr}
Soient $(X,\mathcal{A})$ un espace annelé, $\mathcal{E}$ un
$\mathcal{A}$-Module \emph{localement libre} (de rang fini), $u$ un
endomorphisme de $\mathcal{E}$. Il y a par hypothèse une base $\mathfrak{B}$
de la topologie de $X$ telle que pour tout $V\in\mathfrak{B}$,
$\mathcal{E}|V$ soit isomorphe à $\mathcal{A}^n|V$ (pour un $n$ pouvant varier
avec $V$). Soit $u$ un endomorphisme de $\mathcal{E}$~; pour tout
$V\in\mathfrak{B}$, $u_V$ est donc un endomorphisme du
$\Gamma(V,\mathcal{A})$-module $\Gamma(V,\mathcal{E})$, qui est libre par
hypothèse~; le déterminant $\det u_V$ est donc défini et appartient à
$\Gamma(V,\mathcal{A})$. En outre, si $e_1,\ldots,e_n$ forment une base de
$\Gamma(V,\mathcal{E})$, leurs restrictions à tout ouvert $W\subset V$
forment une base de $\Gamma(W,\mathcal{E})$ sur $\Gamma(W,\mathcal{A})$,
donc $\det u_W$ est la restriction de $\det u_V$ à $W$. Il existe donc une
section et une seule de $\mathcal{A}$ au-dessus de $X$, que nous noterons
$\det u$ et appellerons le \emph{déterminant de $u$}, telle que la restriction
de $\det u$ à tout $V\in\mathfrak{B}$ soit $\det u_V$. Il est clair que pour
tout $x\in X$, on a $(\det u)_x=\det u_x$~; pour deux endomorphismes $u$, $v$
de $\mathcal{E}$, on a
\[
\label{II.6.4.8.1-fr}
\det(u\circ v)=(\det u)(\det v)
\tag{6.4.8.1}
\]
ainsi que
\[
\label{II.6.4.8.2-fr}
\det(1_{\mathcal{E}})=1_{\mathcal{A}}
\tag{6.4.8.2}
\]
et, si le rang de $\mathcal{E}$ est constant (ce qui sera le cas
(\hyperref[0.5.4.1-fr]{0, 5.4.1}) pour $X$ connexe) et égal à $n$
\[
\label{II.6.4.8.3-fr}
\det(s\cdot u)=s^n\det u
\tag{6.4.8.3}
\]
pour tout $s\in\Gamma(X,\mathcal{A})$ (on notera que
$\det(0)=0_{\mathcal{A}}$ si $n\geq1$, mais $\det(0)=1_{\mathcal{A}}$ pour
$n=0$). En outre, pour que $u$ soit un \emph{automorphisme} de
$\mathcal{E}$, il faut et il suffit que $\det u$ soit \emph{inversible} dans
$\Gamma(X,\mathcal{A})$.

Si le rang de $\mathcal{E}$ est constant, on peut définir de la même manière
les fonctions symétriques élémentaires $\sigma_i(u)$, qui sont des éléments
de $\Gamma(X,\mathcal{A})$~; on a encore les relations
(\hyperref[II.6.4.1.3-fr]{6.4.1.3}) à
(\hyperref[II.6.4.1.5-fr]{6.4.1.5}).

